% ============ 9. Conclusion ============
\section{Conclusion}\label{sec:conclusion}

This paper took a working result~---LLM news sentiment predicts next-day stock returns, and beats dictionaries doing it~---and asked the question the result's popularity has outrun: what is the score actually made of? On a frozen corpus of 341 headlines covering 36 U.S.\ mega-caps over 74 trading days, scored under conditions that each vary one ingredient of the measurement and then re-scored in pre-specified staged sessions (interleaved and composition-controlled), the audit returns a five-part answer. The baseline association is marginal, specification-sensitive, and draw-dependent: across fourteen independent re-scorings of the identical headlines (sixteen archived draws including two within-condition means) the coefficient is stable but the $t$-statistic ranges over $[1.47, 2.12]$, so nominal significance depends on the day the pipeline scores ($t=2.10$, BH $q=0.062$, $t=1.82$ with the same-day control on the baseline draw). A 4.75-point named-versus-anonymized gap that initially looked like a firm prior decomposes, under the same-session benchmark, into $+5.79$ points of between-session drift plus a $-1.04$ residual~---and the staged experiments then quantified both layers and bounded the mechanism: the drift is itself mostly a $+3.6$-point \emph{batch-composition} effect (pooled $z=3.4$ across two sittings with inverted run order), and the level effect it masked is real, negative, and two to three points ($-2.82$ interleaved; $-2.05$ under corrected anonymization, direction and magnitude consistent), pointing to the removal of a concrete name rather than a material prior about the identity it carries~---the pre-specified swap control moves levels by essentially nothing ($+0.34$) while the placeholder rewrite moves them by two-plus points, and a sixth-session fictional-name control splits the displacement about evenly between the real name and the rewrite itself. The instrument is reliable but stochastic (ICC $0.919$ pooled, $0.949$ within every sitting; 29.9\% exact replication across four calls), its reliability is stable even while its levels move, and its noise decomposes into a response-level mood shared by each batch's 25 scores (13\% of variance) plus idiosyncratic noise. The predictive loading directionally concentrates in the 15\% of headlines that describe realized price moves~---a classification that survives its own audit, a split that no single firm or day carries, and a difference ($p=0.10$) that remains directional rather than established at this sample size~---and the same-day co-movement of sentiment and returns survives anonymization, placing that channel in the text rather than the model's memory. And coverage itself is selected on realized activity, so the sample over-weights the events where non-semantic continuation and reversal are strongest.

The contribution is not a new predictor. It is a demonstration, on a real corpus with real returns, that the channels a measurement-validity literature has been studying separately~---identity effects, sampling noise, text-level echoes, sample selection, and (newly quantified here) batch context~---can be measured jointly and cheaply, and that doing so changes the interpretation of a headline statistic~---in our own case, it retired a session-confounded reading that the two-session design had produced, and then caught the confound's own confound when the staged experiment attributed most of the ``session drift'' to batch composition. The protocol that fell out of this exercise is deliberately minimal: anonymize once, verify the anonymization, and field the contrast interleaved in one sitting; hold the batch composition fixed across anything you compare; repeat-call once and report the noise floor; flag the text for realized-price language and audit the flag against a written rulebook; test the sample for selection; and pre-register the family before looking at returns. Every component is implementable in an afternoon at the marginal cost of one additional scoring pass~---and the interleaved re-fielding, the one component that costs a full session, is the one this paper needed to correct its own sign.

Two conclusions for the literature. For studies that use LLM sentiment as a return predictor: a significant coefficient on a single-call, named, search-collected corpus is not evidence about text semantics until the echo and selection channels are ruled out; the burden of proof is now demonstrably cheap to meet. For the emerging audit literature: the post-cutoff window is the clean laboratory for identity-channel experiments, and two signatures documented here~---the mild negative de-identification effect (two to three points, replicated across three controlled sittings, and decomposed by the swap and fictional-name controls into a real-name half and a rewrite half) and the batch-composition effect ($+3.6$ points for re-batching alone, robust to inverted run order)~---are precise, cheap targets for cross-model and cross-market replication, with the open question of whether the composition effect's non-linear mechanism generalizes. The questions worth asking have shifted from whether LLMs can read the news to what, exactly, they are reading~---and how much of the market, the firm, and the neighboring headlines is in the score before the text ever arrives.
